\clearpage\begingroup\NoSourceHead\SourceLeaf{541b}
\renewcommand{\thefootnote}{\arabic{footnote})}\setcounter{footnote}{0}
\ArithmeticHead{84}{L}{BOLZA, GAUSS UND DIE VARIATIONSRECHNUNG.}
\begin{center}\bfseries IV. Teil: Vereinzelte kürzere Anwendungen der Variationsrechnung.\end{center}
\begin{center}\textls[100]{1. Fragment: Über das Extremum des Integrals} $\int n(x,y)\,ds$.\\
\small (Handbuch 21, S. 48. Werke, Bd. XI, 1.)\end{center}
\noindent\hspace*{\parindent}\textsc{Gauss} leitet in diesem Fragment die \textsc{Euler}sche Differentialgleichung für
das ebene Variationsproblem
\[
\int n(x,y)\,ds=\mathrm{Min.},\quad ds=\sqrt{dx^2+dy^2}
\]
ab. Er betrachtet $x$ als die unabhängige Variable, sodass das Integral lautet
\[
\int n(x,y)\sqrt{1+p^2}\,dx,\quad p=\frac{dy}{dx}.
\]
\par\noindent\hspace*{\parindent}Als Resultat erhält er nach Einführung des Tangentenwinkels $\varphi$ die
Differentialgleichung in der Form
\begin{equation*}\tag{1}
\frac{\partial n}{\partial y}\cos\varphi-\frac{\partial n}{\partial x}\sin\varphi=n\frac{d\varphi}{ds},
\end{equation*}
d. h. also
\begin{equation*}\tag{2}
\frac{\partial n}{\partial y}\cos\varphi-\frac{\partial n}{\partial x}\sin\varphi=\frac{n}{r},
\end{equation*}
wo $\frac1r$ die Krümmung bedeutet.
\par\noindent\hspace*{\parindent}Eigentümlich ist bei der Ableitung einmal, dass \textsc{Gauss} nicht von der
fertig vorliegenden \textsc{Euler}schen Differentialgleichung für den allgemeinen Typus
\clearpage\SourceLeaf{542}
\ArithmeticHead{85}{R}{VEREINZELTE ANWENDUNGEN DER VARIATIONSRECHNUNG.}
\noindent $\int f(x,y,p)\,dx$ Gebrauch macht, sondern sie für den besonderen Fall durch
Variation des Integrals und Anwendung der \textsc{Lagrange}schen partiellen Integration selbst ableitet, und andererseits die Einführung des Tangentenwinkels $\varphi$.
\par\noindent\hspace*{\parindent}Auf die Differentialgleichung (2) wird man auch unmittelbar geführt,
wenn man das Problem in Parameterdarstellung ansetzt und dafür die \textsc{Weierstrass}sche Form der \textsc{Euler}schen Differentialgleichung hinschreibt.
\par\noindent\hspace*{\parindent}Die Differentialgleichung (2) findet sich in etwas veränderter Form bei
\textsc{Jellett}\footnote{\textit{An elementary Treatise on the Calculus of Variations}, Dublin 1850, S. 140.} wieder, der statt des Tangentenwinkels die Winkel der Kurvennormale mit den beiden Koordinatenachsen einführt. Das Resultat von \textsc{Jellett} hat dann schliesslich W. \textsc{Thomson}\footnote{\textit{Isoperimetrical Problems}, Nature 1894, S. 517; vgl. auch v. \textsc{Schrutka}, \textit{Die ökonomischste Trassenführung für den Fall, dass die Kosten für den laufenden Kilometer mit dem Orte wechseln}, Österreichische
Wochenschrift für den öffentlichen Baudienst, Heft 37 (1911), S. 1—8.} in die einfache Form setzt:
\[
\frac1r=\frac{\partial\log n}{\partial\nu},
\]
wo $\nu$ diejenige Richtung der Kurvennormale bedeutet, die aus der positiven
Tangentenrichtung durch eine positive Drehung um $\frac\pi2$ hervorgeht.
\par\noindent\hspace*{\parindent}Wie aus der Wahl des Buchstabens $n$ hervorgeht, hat \textsc{Gauss} ohne Zweifel
die Anwendung des Problems auf die Bestimmung der Bahn des Lichtes in
einem durchsichtigen Medium mit stetig veränderlichem Brechungsindex $n$ im
Auge gehabt.
\clearpage\endgroup
